% Prajñāpāramitāhṛdaya - Heart Sūtra
% Chinese Critical Edition
% Following Chinese Compositional Priority (Nattier 1992)

\documentclass[11pt,a4paper]{article}

% ============================================================================
% PACKAGES
% ============================================================================

\usepackage[margin=2.5cm,bottom=3cm]{geometry}
\usepackage{fontspec}
\usepackage{xeCJK}
\usepackage{xcolor}
\usepackage{longtable}
\usepackage{array}
\usepackage{enumitem}

% Per-page footnotes
\usepackage[perpage]{footmisc}

% Page style
\usepackage{fancyhdr}
\pagestyle{fancy}
\fancyhf{}
\renewcommand{\headrulewidth}{0.4pt}
\fancyhead[L]{\small Prajñāpāramitāhṛdaya -- Chinese Critical Edition}
\fancyhead[R]{\small 般若波羅蜜多心經}
\fancyfoot[C]{\thepage}

% ============================================================================
% FONTS
% ============================================================================

\setmainfont{Noto Serif}[
  BoldFont = {Noto Serif},
  BoldFeatures = {RawFeature=+axis{wght=700}},
]
\setsansfont{Noto Sans}
\setCJKmainfont{Noto Serif CJK SC}
\setCJKsansfont{Noto Sans CJK SC}

% ============================================================================
% CUSTOM COMMANDS
% ============================================================================

\newcommand{\linenum}[1]{\textbf{\small #1.}}
\newcommand{\pinyin}[1]{{\small\color{gray} #1}}
\newcommand{\var}[1]{\textsuperscript{\color{red}\textbf{#1}}}
\newcommand{\lemma}[1]{\textbf{#1}}
\newcommand{\reading}[1]{#1}
\newcommand{\wit}[1]{\textsc{#1}}
\newcommand{\gloss}[1]{{\small\color{teal}\textit{#1}}}
\newcommand{\fluent}[1]{{\small #1}}

% ============================================================================
% DOCUMENT
% ============================================================================

\begin{document}

% ----------------------------------------------------------------------------
% TITLE PAGE
% ----------------------------------------------------------------------------

\begin{center}
{\Huge\bfseries 般若波羅蜜多心經}\\[0.5em]
{\LARGE Prajñāpāramitāhṛdaya}\\[0.3em]
{\Large Heart Sūtra}\\[1.5em]
{\large Chinese Critical Edition}\\[0.5em]
{\normalsize Based on T251 (Xuanzang, ca. 649--656 CE)}\\[2em]
\end{center}

\hrule
\vspace{1em}

\section*{Witnesses}

\begin{description}[leftmargin=2cm, labelwidth=1.8cm]
\item[\wit{T250}] Dàmíngzhòujīng 大明呪經 (attrib. Kumārajīva, actual date uncertain). Parallel composition preserving three-epithet formula. Contains skandha-characteristic and temporal negation passages absent from T251.
\item[\wit{T251}] Xīnjīng 心經 (Xuanzang, ca. 649--656 CE). Base text. Canonical short recension.
\item[\wit{Fangshan}] Fangshan Stele 房山石經 (661 CE). Oldest dated physical witness.
\item[\wit{T256}] Tang Transliteration (8th--9th c.). Sanskrit in Chinese characters.
\item[\wit{T257}] Dānapāla Translation 施護譯 (1005 CE). Long recension from Sanskrit.
\item[\wit{S.2464}] Dunhuang MS (7th--8th c.). T251-type.
\item[\wit{S.700}] Dunhuang MS (7th--8th c.). T250-type.
\item[\wit{P.2323}] Dunhuang MS (851 CE). Dated colophon.
\end{description}

\vspace{1em}
\noindent\textbf{Note:} T250 and T251 are \emph{parallel compositions} from Large Sūtra material (T223), not stages in a single text's revision history. T250 contains passages absent from T251 (skandha characteristics, temporal negation) and preserves the original three-epithet formula.

\newpage

% ----------------------------------------------------------------------------
% TEXT
% ----------------------------------------------------------------------------

\section*{Text}

\begin{center}
{\Large\bfseries 般若波羅蜜多心經}\\[0.3em]
{\normalsize Bōrě Bōluómìduō Xīnjīng}
\end{center}

\vspace{1em}

\subsection*{Section I: Opening}

\linenum{1} 觀自在\footnote{\lemma{觀自在} \textit{Guānzìzài}] \reading{觀世音} \textit{Guānshìyīn} \wit{T250}, \wit{T252}. The reading 觀自在 supports Xuanzang's Sanskrit \textit{avalokita-īśvara} (``Lord who looks down''), representing his new translation of this name. T250's 觀世音 reflects older \textit{avalokita-svara} (``Perceiver of sounds''). This terminological difference confirms T250 predates Xuanzang's translation conventions.}菩薩，\\
\pinyin{Guānzìzài púsà,}\\[0.3em]
\linenum{2} 行深般若波羅蜜多時，\\
\pinyin{xíng shēn bōrě bōluómìduō shí,}\\[0.3em]
\linenum{3} 照見五蘊\footnote{\lemma{五蘊} \textit{wǔyùn}] \reading{五陰} \textit{wǔyīn} \wit{T250}. T250 uses the older term 五陰 (literally ``five shades/aggregates''), while T251 uses Xuanzang's preferred 五蘊. Both translate Sanskrit \textit{pañca-skandha}.}皆空，\\
\pinyin{zhàojiàn wǔyùn jiē kōng,}\\[0.3em]
\linenum{4} 度一切苦厄\footnote{\lemma{度一切苦厄}] This phrase has no parallel in the Sanskrit text or Large Sūtra source. It appears to be unique to the Chinese Heart Sūtra composition, perhaps added as a liturgical enhancement.}。\\
\pinyin{dù yīqiè kǔ è.}

\vspace{0.3em}
\noindent{\small\color{gray}[\,\textit{T250 inserts skandha-characteristic passage here (see Appendix)}\,]}

\vspace{0.5em}
\gloss{觀自在 (contemplator of self-mastery) 菩薩 (bodhisattva), 行 (practicing) 深 (profound) 般若波羅蜜多 (prajñāpāramitā) 時 (when), 照見 (illuminated and saw) 五蘊 (five aggregates) 皆 (all) 空 (empty), 度 (crossed beyond) 一切 (all) 苦厄 (suffering and distress).}

\vspace{0.3em}
\fluent{When Avalokiteśvara Bodhisattva was practicing the profound prajñāpāramitā, he illuminated the five aggregates and saw that they are all empty, and crossed beyond all suffering and distress.}

\vspace{1em}
\subsection*{Section II: Form and Emptiness}

\linenum{5} 舍利子，\\
\pinyin{Shèlìzǐ,}\\[0.3em]
\linenum{6} 色不異空，空不異色；\\
\pinyin{sè bù yì kōng, kōng bù yì sè;}\\[0.3em]
\linenum{7} 色即是空，空即是色。\\
\pinyin{sè jí shì kōng, kōng jí shì sè.}\\[0.3em]
\linenum{8} 受想行識亦復如是。\\
\pinyin{Shòu xiǎng xíng shí yì fù rúshì.}

\vspace{0.5em}
\gloss{舍利子 (Śāriputra), 色 (form) 不 (not) 異 (different from) 空 (emptiness), 空 (emptiness) 不 (not) 異 (different from) 色 (form); 色 (form) 即是 (is precisely) 空 (emptiness), 空 (emptiness) 即是 (is precisely) 色 (form). 受 (feeling) 想 (perception) 行 (formation) 識 (consciousness) 亦復如是 (are also like this).}

\vspace{0.3em}
\fluent{Śāriputra, form is not different from emptiness, emptiness is not different from form; form is emptiness, emptiness is form. Feeling, perception, formation, and consciousness are also like this.}

\vspace{1em}
\subsection*{Section III: Characteristics of Emptiness}

\linenum{9} 舍利子，是諸法空相，\\
\pinyin{Shèlìzǐ, shì zhūfǎ kōng xiāng,}\\[0.3em]
\linenum{10} 不生不滅，不垢不淨，不增不減。\\
\pinyin{bù shēng bù miè, bù gòu bù jìng, bù zēng bù jiǎn.}

\vspace{0.3em}
\noindent{\small\color{gray}[\,\textit{T250 inserts temporal negation passage here (see Appendix)}\,]}

\vspace{0.5em}
\gloss{舍利子 (Śāriputra), 是 (these) 諸法 (all dharmas) 空相 (emptiness-characteristic): 不生 (not arising) 不滅 (not ceasing), 不垢 (not defiled) 不淨 (not pure), 不增 (not increasing) 不減 (not decreasing).}

\vspace{0.3em}
\fluent{Śāriputra, the emptiness-characteristic of all dharmas is: not arising, not ceasing; not defiled, not pure; not increasing, not decreasing.}

\vspace{1em}
\subsection*{Section IV: Negation Series}

\linenum{11} 是故空中，無色，無受想行識，\\
\pinyin{Shìgù kōng zhōng, wú sè, wú shòu xiǎng xíng shí,}\\[0.3em]
\linenum{12} 無眼耳鼻舌身意，無色聲香味觸法，\\
\pinyin{wú yǎn ěr bí shé shēn yì, wú sè shēng xiāng wèi chù fǎ,}\\[0.3em]
\linenum{13} 無眼界，乃至無意識界。\\
\pinyin{wú yǎn jiè, nǎizhì wú yìshí jiè.}\\[0.3em]
\linenum{14} 無無明，亦無無明盡\footnote{\lemma{盡} \textit{jìn} ``ending''] Key evidence for Chinese priority and Sanskrit back-translation. The Chinese 盡 was translated into Sanskrit as \textit{kṣaya} (``destruction, wearing away'') rather than the standard Prajñāpāramitā term \textit{nirodha} (``cessation''). The Large Sūtra (Pañcaviṃśatisāhasrikā) uses \textit{nirodha}; Kumārajīva's T223 translates this as 盡. A back-translator working from Chinese chose \textit{kṣaya} as a synonym, creating the ``smoking gun'' evidence (Nattier 1992:175--178).}，乃至無老死，亦無老死盡。\\
\pinyin{Wú wúmíng, yì wú wúmíng jìn, nǎizhì wú lǎosǐ, yì wú lǎosǐ jìn.}\\[0.3em]
\linenum{15} 無苦集滅道，無智亦無得\footnote{\lemma{無智亦無得}] Sanskrit adds \textit{na-aprāptiḥ} (``no non-attainment''), absent from T251, Fangshan stele, and Hōryūji Sanskrit MS. This is a later interpolation in the Sanskrit tradition, subsequently translated back into Chinese in T257 as 亦無無得.}，以無所得故。\\
\pinyin{Wú kǔ jí miè dào, wú zhì yì wú dé, yǐ wú suǒ dé gù.}

\vspace{0.5em}
\gloss{是故 (therefore) 空中 (in emptiness): 無 (no) 色 (form), 無 (no) 受 (feeling) 想 (perception) 行 (formation) 識 (consciousness); 無 (no) 眼 (eye) 耳 (ear) 鼻 (nose) 舌 (tongue) 身 (body) 意 (mind), 無 (no) 色 (form) 聲 (sound) 香 (smell) 味 (taste) 觸 (touch) 法 (dharma); 無 (no) 眼界 (eye-element) 乃至 (up to) 無 (no) 意識界 (mind-consciousness element). 無 (no) 無明 (ignorance) 亦 (also) 無 (no) 無明盡 (ending of ignorance), 乃至 (up to) 無 (no) 老死 (old age and death) 亦 (also) 無 (no) 老死盡 (ending of old age and death). 無 (no) 苦 (suffering) 集 (origination) 滅 (cessation) 道 (path); 無 (no) 智 (wisdom) 亦 (also) 無 (no) 得 (attainment), 以 (because) 無所得 (nothing to attain) 故 (therefore).}

\vspace{0.3em}
\fluent{Therefore, in emptiness there is no form, no feeling, perception, formation, or consciousness; no eye, ear, nose, tongue, body, or mind; no form, sound, smell, taste, touch, or mental object; no eye-element, up to and including no mind-consciousness element. No ignorance and no ending of ignorance, up to and including no old age and death and no ending of old age and death. No suffering, origination, cessation, or path; no wisdom and no attainment --- because there is nothing to attain.}

\vspace{1em}
\subsection*{Section V: Bodhisattva's Result}

\linenum{16} 菩提薩埵，依般若波羅蜜多故。\\
\pinyin{Pútísàduǒ, yī bōrě bōluómìduō gù.}\\[0.3em]
\linenum{17} 心無罣礙，無罣礙故。\\
\pinyin{Xīn wú guà'ài, wú guà'ài gù.}\\[0.3em]
\linenum{18} 無有恐怖，遠離顚倒夢想，究竟涅槃。\\
\pinyin{Wú yǒu kǒngbù, yuǎnlí diāndǎo mèngxiǎng, jiūjìng nièpán.}

\vspace{0.5em}
\gloss{菩提薩埵 (bodhisattva), 依 (relying on) 般若波羅蜜多 (prajñāpāramitā) 故 (therefore). 心 (mind) 無 (without) 罣礙 (obstruction), 無 (without) 罣礙 (obstruction) 故 (therefore): 無有 (without) 恐怖 (fear), 遠離 (far removed from) 顚倒 (inverted) 夢想 (dream-thinking), 究竟 (ultimate) 涅槃 (nirvāṇa).}

\vspace{0.3em}
\fluent{Because bodhisattvas rely on prajñāpāramitā, their minds are without obstruction. Without obstruction, they have no fear, are far removed from inverted dream-thinking, and attain ultimate nirvāṇa.}

\vspace{1em}
\subsection*{Section VI: Buddha's Result}

\linenum{19} 三世諸佛，依般若波羅蜜多故，\\
\pinyin{Sānshì zhū fó, yī bōrě bōluómìduō gù,}\\[0.3em]
\linenum{20} 得阿耨多羅三藐三菩提。\\
\pinyin{dé ānòuduōluó sānmiǎo sānpútí.}

\vspace{0.5em}
\gloss{三世 (three times) 諸佛 (all buddhas), 依 (relying on) 般若波羅蜜多 (prajñāpāramitā) 故 (therefore), 得 (attain) 阿耨多羅三藐三菩提 (anuttarā-samyak-saṃbodhi).}

\vspace{0.3em}
\fluent{All buddhas of the three times, relying on prajñāpāramitā, attain unsurpassed, perfect, complete enlightenment.}

\vspace{1em}
\subsection*{Section VII: Mantra Praise}

\linenum{21} 故知般若波羅蜜多，\\
\pinyin{Gùzhī bōrě bōluómìduō,}\\[0.3em]
\linenum{22} 是大神咒\footnote{\lemma{是大神咒，是大明咒，是無上咒，是無等等咒}] T251's four-epithet formula. \wit{T250} has only three: 是大明呪、無上明呪、無等等明呪---\emph{exactly matching} the Large Sūtra (T223). T251 adds 大神咒 and drops 明 from two epithets. The Sanskrit \textit{mahāmantro mahāvidyāmantro 'nuttaramantro 'samasamamantraḥ} follows T251's four-epithet pattern, confirming it was back-translated from T251, not T250.}，是大明咒，是無上咒，是無等等咒，\\
\pinyin{shì dà shén zhòu, shì dà míng zhòu, shì wúshàng zhòu, shì wú děngděng zhòu,}\\[0.3em]
\linenum{23} 能除一切苦，眞實不虛故。\\
\pinyin{néng chú yīqiè kǔ, zhēnshí bù xū gù.}

\vspace{0.5em}
\gloss{故 (therefore) 知 (know) 般若波羅蜜多 (prajñāpāramitā): 是 (is the) 大神咒 (great divine mantra), 是 (is the) 大明咒 (great illumination mantra), 是 (is the) 無上咒 (unsurpassed mantra), 是 (is the) 無等等咒 (unequalled mantra); 能 (able to) 除 (remove) 一切 (all) 苦 (suffering), 眞實 (true) 不虛 (not false) 故 (because).}

\vspace{0.3em}
\fluent{Therefore know that prajñāpāramitā is the great divine mantra, the great illumination mantra, the unsurpassed mantra, the unequalled mantra, which can remove all suffering; it is true, because it is not false.}

\vspace{1em}
\subsection*{Section VIII: Mantra}

\linenum{24} 說般若波羅蜜多咒，卽說咒曰：\\
\pinyin{Shuō bōrě bōluómìduō zhòu, jí shuō zhòu yuē:}\\[0.3em]
\linenum{25} 揭帝揭帝\footnote{\lemma{揭帝} \textit{jiēdì}] Phonetic transcription of Sanskrit \textit{gate}. Minor variants exist in Dunhuang manuscripts. The mantra is transliterated, not translated, in all Chinese witnesses.}，波羅揭帝，波羅僧揭帝，菩提僧莎訶\footnote{Note on \textit{oṃ}: The mantra in T251 begins directly with 揭帝 (\textit{gate}). The \textit{oṃ} found in later Sanskrit and Tibetan witnesses is a Tantric addition (9th--10th century), absent from all early Chinese witnesses, Fangshan stele (661 CE), and Hōryūji Sanskrit MS (7th--8th c.).}。\\
\pinyin{Jiēdì jiēdì, bōluó jiēdì, bōluó sēng jiēdì, pútí sēng suōhē.}

\vspace{0.5em}
\gloss{說 (proclaim) 般若波羅蜜多 (prajñāpāramitā) 咒 (mantra), 卽 (then) 說 (proclaim) 咒 (mantra) 曰 (saying): 揭帝揭帝 (gate gate) 波羅揭帝 (pāragate) 波羅僧揭帝 (pārasaṃgate) 菩提僧莎訶 (bodhi svāhā).}

\vspace{0.3em}
\fluent{Proclaim the prajñāpāramitā mantra, proclaim the mantra that says: Gone, gone, gone beyond, gone completely beyond --- awakening! Svāhā!}

\newpage

% ----------------------------------------------------------------------------
% T250 UNIQUE PASSAGES
% ----------------------------------------------------------------------------

\section*{Appendix: T250 Unique Passages}

T250 contains two substantial passages absent from T251, providing key evidence that T250 is a \emph{parallel composition}, not an earlier draft revised into T251.

\subsection*{Skandha-Characteristic Passage}

Following the negation of the five aggregates, T250 adds:

\begin{quote}
\large
色空故無惱壞相，受空故無受相，想空故無知相，行空故無作相，識空故無覺相。

\normalsize
\pinyin{Sè kōng gù wú nǎohuài xiāng, shòu kōng gù wú shòu xiāng, xiǎng kōng gù wú zhī xiāng, xíng kōng gù wú zuò xiāng, shí kōng gù wú jué xiāng.}

\textit{``Because form is empty there is no characteristic of affliction; because feeling is empty there is no characteristic of receiving; because perception is empty there is no characteristic of knowing; because volition is empty there is no characteristic of acting; because consciousness is empty there is no characteristic of awareness.''}
\end{quote}

\subsection*{Temporal Negation Passage}

T250 also includes:

\begin{quote}
\large
是空法，非過去，非未來，非現在。

\normalsize
\pinyin{Shì kōng fǎ, fēi guòqù, fēi wèilái, fēi xiànzài.}

\textit{``This emptiness-dharma is not past, not future, not present.''}
\end{quote}

\noindent These passages represent substantive doctrinal content. A reviser creating T251 from T250 would have no reason to delete them. Their absence from T251 indicates T251 was composed independently, drawing on the same Large Sūtra source material but making different compositional choices.

\newpage

\section*{Note on Prior Editions}

Standard references for the Chinese Heart Sūtra include the Taishō edition (T8.848--849 for T251; T8.847 for T250) and the CBETA digital corpus. Attwood (2017b, 2021a, 2024) has provided the most sustained modern analysis of the Chinese text, including:

\begin{itemize}
\item Systematic comparison of Kumārajīva-era and Xuanzang-era translation conventions
\item Identification of the epithet pattern (3 vs.\ 4 epithets) as key evidence for textual relationships
\item Analysis of 明呪 \textit{míngzhòu} / 神咒 \textit{shénzhòu} terminology
\item A revised edition of the Chinese text in Attwood (2024)
\end{itemize}

The two editions differ in editorial philosophy. Attwood (2024) presents a \emph{revised} edition: he proposes specific emendations to the received Taishō text, replacing 無智亦無得 with Xuanzang's rendering 無得無現觀 (from T220), restoring 明呪 in the epithet passage, correcting 僧莎訶 to 薩婆訶 in the dhāraṇī, and repunctuating 以無所得故 to open a new paragraph. These changes aim to produce a text that can be parsed and translated correctly, and are supported by careful argumentation.

This edition takes a different approach: it is a \emph{critical} edition in the traditional philological sense. We present the received T251 text without emendation, recording variants in the apparatus. Where Attwood proposes changing the text, we note the problem in a footnote and let readers evaluate the evidence. Our further contributions beyond Attwood's work:

\begin{itemize}
\item We present T251 and T250 as \emph{parallel compositions}, with T250 readings in the apparatus rather than treating either text as a revision of the other
\item We include the Fangshan stele (661 CE), Dunhuang manuscripts (S.2464, S.700, P.2323), and Dānapāla's T257 as formal witnesses
\item We mark the locations where T250 preserves passages absent from T251 and present them in full
\item Pinyin romanization and word-by-word glosses are provided throughout
\end{itemize}

\vspace{2em}
\hrule
\vspace{0.5em}
\begin{center}
\textit{Chinese Critical Edition} $\cdot$ Prajñāpāramitāhṛdaya $\cdot$ Following Nattier (1992)
\end{center}

\end{document}
